\chapter{Полный гессиан исправленного минимума и относительные семейные моды}

\section{Цель проверки}

Канонический подъём доказал существование rank-$21$ минимума, но предъявил
только факторизованный свидетель. Теперь разрешим трём неэквивалентным
физическим рёбрам независимые аффинные поля
\begin{equation}
 X_u,X_d,X_e\in
 E_{\rm aff}=\operatorname{Hom}(\mathbb C^4,\operatorname{im}P_3)
 \label{eq:v7-three-independent-affine-edges}
\end{equation}
и вычислим весь вещественный гессиан размерности
\begin{equation}
 3\cdot2\cdot3\cdot4=72.
\end{equation}
Проверка должна отличить настоящую устойчивость от плоского семейства
неэквивалентных минимумов.

\section{Точная блочная редукция}

Зафиксируем нормированное значение общего хиггсовского дублета. Три
эквивариантных ребра $T_u,T_d,T_e:H_L^8\to H_R^7$ имеют ортогональные
правые концы и кратности
\begin{equation}
 m_u=3,
 \qquad
 m_d=3,
 \qquad
 m_e=1.
\end{equation}
Поднятый оператор равен
\begin{equation}
 Z(X_u,X_d,X_e)
 =X_u\otimes T_u+X_d\otimes T_d+X_e\otimes T_e.
 \label{eq:v7-three-edge-lifted-operator}
\end{equation}
Ортогональность калибровочных концов уничтожает смешанные следы между
разными рёбрами. Поэтому функционал предыдущей главы точно редуцируется к
\begin{equation}
 \begin{split}
 45\mathcal S_{\rm edge}
 =3+\sum_{a\in\{u,d,e\}}m_a\{&
 \|X_a^\dagger X_a-P_3\|_F^2\\
 &+\|X_aX_a^\dagger-I_3\|_F^2\}.
 \end{split}
 \label{eq:v7-relative-edge-action}
\end{equation}
Постоянное слагаемое три является отсутствующим правым нейтринным каналом.

\section{Полное вакуумное многообразие}

Каждое слагаемое минимально тогда и только тогда, когда
\begin{equation}
 X_aX_a^\dagger=I_3,
 \qquad
 X_a^\dagger X_a=P_3.
 \label{eq:v7-edge-coisometry-conditions}
\end{equation}
Поскольку $V$ является фиксированной коизометрией с $V^\dagger V=P_3$,
все решения имеют вид
\begin{equation}
 X_a=U_aV,
 \qquad U_a\in U(3).
\end{equation}
Следовательно,
\begin{equation}
 \mathcal M_{\rm vac}=U(3)_u\times U(3)_d\times U(3)_e,
 \qquad
 \dim_{\mathbb R}\mathcal M_{\rm vac}=27.
 \label{eq:v7-relative-edge-vacuum-manifold}
\end{equation}
Функционал фиксирует сингулярные числа и нормы трёх рёбер, но не их
относительные семейные кадры.

\begin{theorem}[Неизолированность исправленного минимума]
Внутри полного исправленного модуля кратностей все глобальные минимумы
образуют по меньшей мере многообразие $U(3)^3$. Минимум $V\otimes Y_\star$
не является изолированным представителем, а только точкой этого
многообразия.
\end{theorem}

\section{Полный вещественный гессиан}

Прямая конечная разность по всем семидесяти двум координатам даёт сырой
координатный спектр
\begin{equation}
 \Spec\Hess_{\rm raw}=
 \left\{
 0^{(27)},
 \left(\frac4{45}\right)^{(6)},
 \left(\frac4{15}\right)^{(12)},
 \left(\frac{16}{45}\right)^{(9)},
 \left(\frac{16}{15}\right)^{(18)}
 \right\}.
 \label{eq:v7-relative-edge-raw-hessian}
\end{equation}
Различие ненулевых уровней между $u,d$ и $e$ происходит только из цветовых
кратностей $3,3,1$. Кинетическая trace-метрика имеет те же веса. Поэтому
физический обобщённый спектр равен
\begin{equation}
 \Spec(G^{-1}\Hess)=
 \left\{
 0^{(27)},
 \left(\frac4{45}\right)^{(18)},
 \left(\frac{16}{45}\right)^{(27)}
 \right\}.
 \label{eq:v7-relative-edge-generalized-hessian}
\end{equation}
Отрицательных направлений нет; сорок пять поперечных направлений массивны,
а двадцать семь касательных направлений точно плоски.

\section{Какие нули могут быть калибровочными}

Текущая семейная структура имеет дискретное $S_4$, а не непрерывную
калибровочную группу $U(3)^3$. Поэтому автоматически удалить двадцать семь
нулей как БРСТ-точные нельзя.

Даже если в максимально благоприятном для кандидата сценарии объявить одно
общее диагональное $U(3)$ калибровочным, остаётся относительный фактор
\begin{equation}
 \frac{U(3)_u\times U(3)_d\times U(3)_e}{U(3)_{\rm diag}},
 \qquad
 \dim_{\mathbb R}=27-9=18.
 \label{eq:v7-relative-unitary-moduli}
\end{equation}
Если Real-условие сокращает унитарные кадры до вещественных ортогональных,
то всё равно остаётся
\begin{equation}
 \dim\frac{O(3)^3}{O(3)_{\rm diag}}=9-3=6
 \label{eq:v7-relative-real-moduli}
\end{equation}
относительных нулевых направлений.

Эти моды являются именно взаимными ориентациями рёбер $u,d,e$. Они
совпадают по типу с данными, из которых впоследствии читаются семейные
матрицы смешивания, но текущий функционал не выбирает их значения.

\begin{nogocriterion}[Запрет чтения CKM/PMNS из плоского многообразия]
Наличие относительных матриц $U_uU_d^\dagger$ и $U_eU_\nu^\dagger$ не
является их выводом. Пока действие постоянно вдоль соответствующих орбит,
любая конкретная матрица смешивания является выбором точки вакуумного
многообразия.
\end{nogocriterion}

\section{Вердикт и следующий различающий тест}

Исправленный Hodge-родитель проходит запуск и поперечную устойчивость, но не
единственность endpoint. Нужный следующий объект должен связывать
неэквивалентные рёбра без независимых портальных коэффициентов.

Обычный линейный интертвинер между $u,d,e$ запрещён их бимодульными метками.
Поэтому остаётся один внутренний маршрут: проверить, создаёт ли уже
допустимая вторая степень, общий хиггсовский дублет или оператор Вайнберга
канонический перекрёстный инвариант, чувствительный к относительным
$U_aU_b^\dagger$. Если полный junk-фактор снова делает функционал блочно
диагональным, смешивание не выводится этой архитектурой.

\begin{protocol}
\begin{enumerate}
 \item полный исправленный вещественный гессиан: \textbf{вычислен};
 \item отрицательных мод в минимуме: \textbf{$0$};
 \item положительных мод: \textbf{$45$};
 \item точных нулевых мод: \textbf{$27$};
 \item вакуумное многообразие: \textbf{$U(3)^3$};
 \item относительных мод после общего $U(3)$: \textbf{$18$};
 \item относительных мод в вещественном пределе: \textbf{$6$};
 \item фиксация сингулярных чисел рёбер: \textbf{получена};
 \item фиксация относительных семейных кадров: \textbf{не получена};
 \item изолированный endpoint: \textbf{не получен};
 \item следующий гейт: \textbf{перекрёстная кривизна второй степени между
 физическими рёбрами}.
\end{enumerate}
\end{protocol}

Машинный сертификат сохраняется в
\path{s2t_v7_corrected_vacuum_relative_edge_hessian_gate_results.json}.